\documentclass[11pt]{article}
\usepackage[margin=1in]{geometry}
\usepackage{amsmath,amssymb,booktabs,graphicx,hyperref,xcolor,framed}
\hypersetup{colorlinks=true,linkcolor=blue!50!black,urlcolor=blue!50!black}
\newcommand{\ActA}{18}
\newcommand{\NptsA}{240}
\newcommand{\GaussA}{1.2\times10^{-8}}
\newcommand{\GaussHiA}{9.6\times10^{-8}}
\newcommand{\BianA}{7.6\times10^{-13}}
\newcommand{\BianHiA}{1.5\times10^{-10}}
\newcommand{\AsymA}{2.9\times10^{-10}}
\newcommand{\SplitA}{1.2\times10^{-15}}
\newcommand{\RconnVsGaussA}{4.8\times10^{-10}}
\newcommand{\EigMinA}{3.4\times10^{-4}}
\newcommand{\EigMaxA}{1.31}
\newcommand{\CondA}{4849}
\newcommand{\LegacyA}{1.74\times10^{3}}
\newcommand{\RmedA}{24.1}
\newcommand{\RloA}{-1027}
\newcommand{\RhiA}{6666}
\newcommand{\AbsRA}{307}
\newcommand{\KmedA}{3.8\times10^{-5}}
\newcommand{\KloA}{-0.007}
\newcommand{\KhiA}{0.022}
\newcommand{\KNegA}{45}
\newcommand{\GeoMsA}{0.60}
\newcommand{\GeoMeA}{1.90}
\newcommand{\GeoLenA}{6.263}
\newcommand{\LineLenA}{6.623}
\newcommand{\GeoNormalA}{94}
\newcommand{\LineNormalA}{3}
\newcommand{\GeoHawkA}{0.108}
\newcommand{\LineHawkA}{0.049}
\newcommand{\GeoEinA}{0.227}
\newcommand{\LineEinA}{0.101}
\newcommand{\GeoSpeedA}{2\times10^{-4}}
\newcommand{\RDustA}{-138.4}
\newcommand{\ShCoDustA}{$<$1}
\newcommand{\ShGrDustA}{14}
\newcommand{\ShQuDustA}{54}
\newcommand{\ShHoDustA}{19}
\newcommand{\ShPlDustA}{13}
\newcommand{\RRadiationA}{43.3}
\newcommand{\ShCoRadiationA}{$<$1}
\newcommand{\ShGrRadiationA}{12}
\newcommand{\ShQuRadiationA}{42}
\newcommand{\ShHoRadiationA}{28}
\newcommand{\ShPlRadiationA}{18}
\newcommand{\RVacuumA}{501.8}
\newcommand{\ShCoVacuumA}{$<$1}
\newcommand{\ShGrVacuumA}{13}
\newcommand{\ShQuVacuumA}{50}
\newcommand{\ShHoVacuumA}{18}
\newcommand{\ShPlVacuumA}{20}
\newcommand{\RBlackholeA}{-15.9}
\newcommand{\ShCoBlackholeA}{$<$1}
\newcommand{\ShGrBlackholeA}{27}
\newcommand{\ShQuBlackholeA}{21}
\newcommand{\ShHoBlackholeA}{32}
\newcommand{\ShPlBlackholeA}{19}
\newcommand{\RWormholeA}{363.4}
\newcommand{\ShCoWormholeA}{$<$1}
\newcommand{\ShGrWormholeA}{16}
\newcommand{\ShQuWormholeA}{37}
\newcommand{\ShHoWormholeA}{26}
\newcommand{\ShPlWormholeA}{21}
\newcommand{\RHawkingA}{-8.6}
\newcommand{\ShCoHawkingA}{$<$1}
\newcommand{\ShGrHawkingA}{25}
\newcommand{\ShQuHawkingA}{29}
\newcommand{\ShHoHawkingA}{28}
\newcommand{\ShPlHawkingA}{17}
\newcommand{\BerryA}{8\times10^{-16}}
\newcommand{\ActB}{18}
\newcommand{\NptsB}{240}
\newcommand{\GaussB}{1.6\times10^{-8}}
\newcommand{\GaussHiB}{1.7\times10^{-7}}
\newcommand{\BianB}{1.1\times10^{-12}}
\newcommand{\BianHiB}{8.2\times10^{-11}}
\newcommand{\AsymB}{4.2\times10^{-10}}
\newcommand{\SplitB}{1.2\times10^{-15}}
\newcommand{\RconnVsGaussB}{4.3\times10^{-10}}
\newcommand{\EigMinB}{2.2\times10^{-4}}
\newcommand{\EigMaxB}{1.29}
\newcommand{\CondB}{7562}
\newcommand{\LegacyB}{2.62\times10^{3}}
\newcommand{\RmedB}{-27.9}
\newcommand{\RloB}{-5496}
\newcommand{\RhiB}{3773}
\newcommand{\AbsRB}{524}
\newcommand{\KmedB}{7.8\times10^{-5}}
\newcommand{\KloB}{-0.006}
\newcommand{\KhiB}{0.025}
\newcommand{\KNegB}{41}
\newcommand{\GeoMsB}{0.60}
\newcommand{\GeoMeB}{1.90}
\newcommand{\GeoLenB}{6.205}
\newcommand{\LineLenB}{6.560}
\newcommand{\GeoNormalB}{97}
\newcommand{\LineNormalB}{5}
\newcommand{\GeoHawkB}{0.118}
\newcommand{\LineHawkB}{0.041}
\newcommand{\GeoEinB}{0.415}
\newcommand{\LineEinB}{0.084}
\newcommand{\GeoSpeedB}{6\times10^{-5}}
\newcommand{\RDustB}{-384.8}
\newcommand{\ShCoDustB}{$<$1}
\newcommand{\ShGrDustB}{14}
\newcommand{\ShQuDustB}{42}
\newcommand{\ShHoDustB}{26}
\newcommand{\ShPlDustB}{18}
\newcommand{\RRadiationB}{-143.4}
\newcommand{\ShCoRadiationB}{$<$1}
\newcommand{\ShGrRadiationB}{13}
\newcommand{\ShQuRadiationB}{38}
\newcommand{\ShHoRadiationB}{30}
\newcommand{\ShPlRadiationB}{18}
\newcommand{\RVacuumB}{-304.5}
\newcommand{\ShCoVacuumB}{$<$1}
\newcommand{\ShGrVacuumB}{8}
\newcommand{\ShQuVacuumB}{37}
\newcommand{\ShHoVacuumB}{31}
\newcommand{\ShPlVacuumB}{24}
\newcommand{\RBlackholeB}{-24.6}
\newcommand{\ShCoBlackholeB}{$<$1}
\newcommand{\ShGrBlackholeB}{22}
\newcommand{\ShQuBlackholeB}{21}
\newcommand{\ShHoBlackholeB}{34}
\newcommand{\ShPlBlackholeB}{22}
\newcommand{\RWormholeB}{519.8}
\newcommand{\ShCoWormholeB}{$<$1}
\newcommand{\ShGrWormholeB}{12}
\newcommand{\ShQuWormholeB}{35}
\newcommand{\ShHoWormholeB}{23}
\newcommand{\ShPlWormholeB}{29}
\newcommand{\RHawkingB}{-3.2}
\newcommand{\ShCoHawkingB}{$<$1}
\newcommand{\ShGrHawkingB}{28}
\newcommand{\ShQuHawkingB}{17}
\newcommand{\ShHoHawkingB}{34}
\newcommand{\ShPlHawkingB}{21}
\newcommand{\BerryB}{7\times10^{-16}}


\title{\textbf{Quantum Mechanics as Geometry:\\ A Geodesic Unification with Gravitation}\\[3pt]
\large Version 2 (November 2026): the principle, realised and measured on a trained M-W manifold}
\author{Rahul Modak$^{1}$ and Rahul Walawalkar$^{1,2}$\\[2pt]
\small $^1$Bayesian Cybersecurity Pvt.\ Ltd., Mumbai, India \quad $^2$NETRA and Caret Capital\\
\small Code: \url{https://github.com/RahulModak74/QUANTUM_GRAVITY_WITH_MW}\\ \small (folder \texttt{Nov-2026\_tests/EXPLORATORY\_PHYSICS})}
\date{}

\begin{document}
\maketitle
\begin{framed}
\small\textbf{Disclaimer.} This is a conceptual note with computational support. We do not claim to
have derived quantum mechanics from geometry or to have solved quantum gravity. Version 1 proposed
a principle. Version 2 tests what that principle implies on one concrete object, the Nov-2026
M-W manifold trained on synthetic GR+QM data, and reports what holds there.
\end{framed}

\begin{abstract}
Einstein showed that gravity is the geometry of spacetime. Version 1 proposed an equivalence
principle in which gravitational and quantum behaviour are both geodesic motion on one manifold
$\mathcal M$, derived from the action
$\int(g_{ij}\dot z^i\dot z^j+\hbar A_i\dot z^i)\,d\lambda$. Version 2 evaluates every term of that
action on a trained M-W manifold. (i) One metric carries both theories: $g$ is a sum of a
gravity part and a quantum part, and so is the curvature, through an Einstein-form equation
$G_{ab}=\kappa(T^{\rm grav}+T^{\rm quant}+\dots)_{ab}$ that holds to $\GaussA$. (ii) Geodesics are
physical: the geodesic between two black holes is a continuous family of black holes. (iii) The
quantum connection built from the decoded wavefunction is pure gauge ($|F|/|\partial A|<\BerryA$),
so the $\hbar A$ term contributes only an endpoint phase. The classical and quantum limits of the
action therefore share the same geodesics on this manifold.
\end{abstract}

\section{The principle (unchanged from Version 1)}
\begin{quote}
All physical systems satisfying both gravitational and quantum constraints evolve as geodesic
flows on a single underlying manifold $\mathcal M$.
\end{quote}
GR identifies gravity with geometry: free particles follow geodesics
$\ddot x^\mu+\Gamma^\mu{}_{\nu\rho}\dot x^\nu\dot x^\rho=0$, and curvature is sourced by
$T_{\mu\nu}$. QM is usually stated algebraically, but its dynamics also come from an extremal
principle (the path integral), and curvature in state space has observable effects (Berry phase).
Version 1 proposed one variational principle for both:
\begin{equation}
\delta\!\int_{\gamma\subset\mathcal M}\bigl(g_{ij}(z)\dot z^i\dot z^j+\hbar A_i(z)\dot z^i\bigr)\,d\lambda=0.
\label{eq:action}
\end{equation}
Version 1 left $\mathcal M$ abstract. Here $\mathcal M$ is a concrete manifold, so every term of
(\ref{eq:action}) can be evaluated.

\section{The manifold}
$\mathcal M$ is the decoder image of the Nov-2026 M-W VAE (checkpoint \texttt{A\_full\_prior}),
trained on 30{,}000 joint GR+QM states from six regimes, with GR and QM constraints as Bayesian
priors. The decoder $f$ maps the \ActA{} active latent axes into 30 standardised columns. The
metric is the pullback $g_{ab}=\langle\partial_af,\partial_bf\rangle$, computed exactly together
with its connection and curvature.

\section{One metric, two theories}
The inner product is a sum over data columns, so the metric decomposes exactly:
\begin{equation}
g_{ab}=g^{\rm grav}_{ab}+g^{\rm quant}_{ab}+g^{\rm hor}_{ab}+g^{\rm coord}_{ab}+g^{\rm Planck}_{ab},\qquad
g^{(s)}_{ab}=\sum_{k\in s}\partial_af^k\,\partial_bf^k.
\end{equation}
The distance between two states is measured jointly by how their gravitational quantities
($g_{\mu\nu}$, $R$, $T_{\mu\nu}$) and their quantum quantities ($\psi$, $\Delta x\Delta t$) differ.
Curvature splits the same way through the Gauss equation, giving an Einstein-form field equation
with additive sources,
\begin{equation}
G_{ab}=\kappa\bigl(T^{\rm grav}+T^{\rm quant}+T^{\rm hor}+T^{\rm coord}+T^{\rm Planck}\bigr)_{ab},
\qquad \nabla^aT_{ab}=0,
\end{equation}
which holds on the trained manifold to $\GaussA$ (Gauss) and $\BianA$ (Bianchi); details in
\emph{Field Equations of the M-W Manifold, Version 2}.

\begin{table}[h]\centering\small
\begin{tabular}{@{}lccccc@{}}\toprule
Regime & gravity & quantum & horizon & Planck & coord.\\\midrule
dust & \ShGrDustA & \ShQuDustA & \ShHoDustA & \ShPlDustA & \ShCoDustA\\
radiation & \ShGrRadiationA & \ShQuRadiationA & \ShHoRadiationA & \ShPlRadiationA & \ShCoRadiationA\\
vacuum & \ShGrVacuumA & \ShQuVacuumA & \ShHoVacuumA & \ShPlVacuumA & \ShCoVacuumA\\
black hole & \ShGrBlackholeA & \ShQuBlackholeA & \ShHoBlackholeA & \ShPlBlackholeA & \ShCoBlackholeA\\
wormhole & \ShGrWormholeA & \ShQuWormholeA & \ShHoWormholeA & \ShPlWormholeA & \ShCoWormholeA\\
Hawking & \ShGrHawkingA & \ShQuHawkingA & \ShHoHawkingA & \ShPlHawkingA & \ShCoHawkingA\\\bottomrule
\end{tabular}
\caption{Share (\%) of the scalar curvature $|R|$ contributed by each column group, mean over 40
unseen test states per regime. Every regime is curved by gravitational and quantum content
together.}
\label{tab:share}
\end{table}

Table~\ref{tab:share} gives the measured version of the principle's central claim: no regime is
curved by gravity alone or by quantum content alone. In the matter regimes the quantum columns
carry the largest share of curvature. In the black-hole and Hawking regimes the gravity and horizon
columns take a larger share than in the matter regimes, while quantum content still contributes.

\section{Geodesics are physical}
The geodesic between a black hole of mass $\GeoMsA$ and one of mass $\GeoMeA$ was found by
minimising path energy. $\GeoNormalA\%$ of its acceleration is normal to the manifold (straight
latent line: $\LineNormalA\%$), as required of a geodesic. Every point along it decodes to a black
hole, with mass increasing uniformly along the path, and the Hawking relation holds along it to
$|8\pi MT_H-1|=\GeoHawkA$. Geodesic motion on $\mathcal M$ moves continuously through physically
meaningful states, as the principle requires.

\section{The quantum connection is pure gauge here}
The decoded wavefunction $\psi=\psi_{\rm re}+i\psi_{\rm im}$ defines the natural $U(1)$ connection
\begin{equation}
A_a=\frac{\operatorname{Im}(\psi^*\partial_a\psi)}{|\psi|^2}=\partial_a\theta,\qquad F_{ab}=\partial_aA_b-\partial_bA_a.
\end{equation}
Evaluated with exact derivatives at 600 test states, $|F|/|\partial A|<\BerryA$ for both checkpoints:
$F=0$, so $A$ is pure gauge, as it must be for a single complex amplitude away from its zeros. The
term $\hbar\!\int\!A_i\dot z^i d\lambda=\hbar\,\Delta\theta$ depends only on the endpoints and
does not change the Euler--Lagrange equations of (\ref{eq:action}). On this manifold the classical
($\hbar\to0$) and quantum regimes of the action share the same geodesics, and the quantum
content enters through the metric (Section~3) rather than through a separate force. A non-zero
Berry curvature needs a multi-component state, i.e.\ more than one complex amplitude per row; we
list this as the next construction.

\section{Uncertainty and non-commutativity}
Version 1 related $[\hat x,\hat p]=i\hbar$ to $[\nabla_i,\nabla_j]V^k=R^k{}_{lij}V^l$. On the
trained manifold the curvature is computed exactly. Sectional curvatures of coordinate planes are
small and of both signs (median $\KmedA$, 5--95\% $[\KloA,\KhiA]$), so covariant derivatives fail
to commute by a measurable, direction-dependent amount. A quantitative link between this and
$\hbar$ is not established here and remains a conjecture.

\section{Conclusion}
Einstein replaced gravitational force with geometry. On a concrete manifold that holds GR and QM
content jointly, both theories now enter one metric and one Einstein-form equation as additive
sources. Geodesics run through physically meaningful states, and the quantum phase term of the
proposed action is pure gauge. This supports, on one learned object, the reading that the
incompatibility between the two theories is a mismatch of descriptions rather than of physics.
It is offered as a working picture to test further, not as a derivation.

\paragraph{Reproduce.} In \texttt{EXPLORATORY\_PHYSICS/code}:
\texttt{mw\_exact\_geometry.py}, \texttt{mw\_berry\_check.py}. Code implementation was assisted by
AI coding tools; modelling choices and claims are the authors'.
\end{document}
